% Evans second edition, section 3.3.
% English source: printed p. 110; PDF page 127 (1-based).
% Chinese section starts: checked against the chapter-opening list, PDF p. 79.
% Status: editorial study draft with detailed proofs and mathematical notes; no independent review.
\Needspace{65mm}
\section{Introduction to Hamilton--Jacobi Equations}
\label{sec:3-3}

\subsection{Variations and Hamilton's ODE (Evans \S3.3.1)}
For a path $X:[0,t]\to\R^n$ with fixed endpoints consider
\[
 \mathcal A[X]=\int_0^t L(\dot X(s),X(s))\,ds.
\]
Suppose a smooth minimizer exists. For every smooth variation $V$
vanishing at both endpoints, differentiate the action at $X+\varepsilon V$:
\[
 0=\int_0^t[L_v(\dot X,X)\cdot\dot V+L_x(\dot X,X)\cdot V]\,ds.
\]
Integration by parts, with the endpoint terms zero, gives
$\frac d{ds}L_v(\dot X,X)=L_x(\dot X,X)$ by the fundamental lemma
of the calculus of variations. These are the Euler--Lagrange equations.
They are necessary conditions; their derivation does not prove a
minimizer exists or that every stationary path minimizes.

Assume now the velocity--momentum relation $P=L_v(v,x)$ can be inverted,
and define
$H(P,x)=P\cdot v(P,x)-L(v(P,x),x)$.
In differentiating this identity the terms containing derivatives of
$v$ cancel because $P=L_v$. Thus $H_p=v$ and $H_x=-L_x$.
The Euler--Lagrange equations become
\[
 \dot X=H_p(P,X),\qquad \dot P=-H_x(P,X).
\]
For $L(v,x)=|v|^2/2-V(x)$ this gives $P=v$,
$H=|P|^2/2+V(x)$ and $\ddot X=-DV(X)$.
For an autonomous Hamiltonian,
$\frac d{ds}H=H_p\cdot(-H_x)+H_x\cdot H_p=0$.
\Needspace{45mm}
\PDEAnnotation{两个来源得到同一组 ODE}{变分法从作用量的驻值推出 Hamilton 方程；特征线法从 Hamilton--Jacobi PDE 推出同一方程。前者给出力学意义，后者解释它如何搬运函数与梯度。}

\subsection{Convex duality and Hopf--Lax (Evans \S3.3.2)}
\label{sec:3-hopflax}
A function $A$ is superlinear if $A(v)/|v|\to+\infty$ as $|v|\to\infty$.
For a finite convex function define
\[
 A^*(p)=\sup_v\{p\cdot v-A(v)\}.
\]
Superlinearity of $A$ makes the objective tend to $-\infty$,
so the maximum is attained and $A^*$ is finite.
It is convex as a supremum of affine functions.
For every $R>0$,
\[
 A^*(p)\geq R|p|-\max_{|v|=R}A(v);
\]
hence $A^*$ is also superlinear when $A$ is finite everywhere.

\begin{theorem}[Convex duality; source Theorem 3]
For finite convex superlinear $A$, $A^{**}=A$ and
\[
 A(v)+A^*(p)\geq p\cdot v.
\]
Equality holds precisely when $p$ is a supporting slope of $A$ at $v$.
For differentiable $A$, this means $p=DA(v)$.
\end{theorem}
\begin{expandedproof}
The inequality is the definition of the supremum.
A supporting slope satisfies
$A(w)\geq A(v)+p\cdot(w-v)$ for every $w$.
Rearranging shows the supremum defining $A^*(p)$ is attained at $v$,
giving equality. Conversely equality rearranges to that support inequality.
Always $A^{**}(v)\leq A(v)$. Choose a supporting slope at $v$;
then the equality just established gives the reverse inequality.
Existence of supporting hyperplanes for finite convex functions is
the invoked convex-analysis prerequisite.
For a differentiable convex function the tangent-plane inequality
provides this slope directly as $DA(v)$.
\end{expandedproof}
The supporting-plane argument for $A^{**}=A$ also holds for every
finite convex $A$ without superlinearity; in that case $A^*$ may take
the value $+\infty$. Subdifferentials of such extended-valued functions
are taken only at points where the function is finite.
\Needspace{45mm}
\PDEAnnotation{对偶变换中的极大值}{求 $v\cdot p-H(p)$ 的驻点只是候选；凸性让该目标关于 $p$ 凹，超线性增长使它在无穷远趋负无穷，才保证极大值确实取得。}

\PDEBodyDiagram{3.2}{Legendre duality is a supporting-line construction.}
{For $H(p)=p^2/2$ and $v=1$, maximizing $vp-H(p)$ finds the tangent point $p=1$. The tangent's intercept is $-L(1)=-1/2$. Convexity makes the whole graph lie above this line, which is exactly the Fenchel inequality.}
{ch03-legendre-support}

For the rest of this subsection assume $H$ is finite, convex, smooth
and superlinear on $\R^n$, set $L=H^*$, and let $g$ be globally
Lipschitz with constant $K$. Then $L$ is finite, convex and superlinear,
and $H=L^*$.

\Needspace{45mm}
\PDEAnnotation{超线性排除逃向远方的起点}{$g$ 的下降至多是线性的，而 $tL((x-y)/t)$ 在位移很大时增长更快。这不仅给出最小值存在，还把最优起点限制在 $|x-y|\leq Rt$ 内。}
\paragraph{Why a straight path gives the formula (source Theorem 4).}
Since $L$ no longer depends on position, Jensen's inequality gives
for any absolutely continuous path with $X(0)=y$, $X(t)=x$,
\[
 \int_0^tL(\dot X(s))\,ds
 \geq tL\left(\frac1t\int_0^t\dot X(s)\,ds\right)
 =tL((x-y)/t).
\]
The straight path $X(s)=y+s(x-y)/t$ attains equality.
Minimizing over its starting point gives
\[
 \boxed{u(x,t)=\min_{y\in\R^n}
             \{tL((x-y)/t)+g(y)\}.}
 \tag{HL}
\]
The minimum exists: with $v=(x-y)/t$, its objective is at least
$g(x)+t[L(v)-K|v|]$, coercive in $v$.
Comparing with $y=x$ also proves that every minimizer satisfies
$L(v)-K|v|\leq L(0)$, hence $|v|\leq R_K$ for a fixed finite $R_K$.
This is the first finite-range localization of the minimization.
\Needspace{45mm}
\PDEAnnotation{直线最优依赖哪些条件}{这里用到 $L$ 的凸性和无位置依赖。若 $L=L(v,x)$，不同路径经过不同势场，不能只用端点位移除以时间就替代整条作用量。}


\PDEMarginDiagram{3.M}{把路径在中途切开}
{中间点 $z$ 也要优化。两段代价相加后再取下确界，才得到动态规划恒等式。}
{ch03-path-split-margin}

\begin{lemma}[Dynamic programming; source Lemma 1]\label{lem:3-dpp}
For $0<s<t$,
\[
 u(x,t)=\min_z\{(t-s)L((x-z)/(t-s))+u(z,s)\}.
 \tag{DP}
\]
\end{lemma}
\begin{expandedproof}
For any $z$ and a minimizing starting point $y$ for $u(z,s)$,
convexity applied to
$(x-y)/t=[(t-s)/t](x-z)/(t-s)+(s/t)(z-y)/s$
gives
$u(x,t)\leq(t-s)L((x-z)/(t-s))+u(z,s)$.
This proves one inequality for the infimum over $z$.
Conversely take a minimizer $y$ in (HL) and choose
$z=y+(s/t)(x-y)$. Both subinterval velocities then equal $(x-y)/t$.
Use this same $y$ as a competitor for $u(z,s)$ to obtain the opposite
inequality. Coercivity ensures the intermediate minimum exists.
\end{expandedproof}

\begin{lemma}[Lipschitz bounds and initial values; source Lemma 2]
The function in (HL) is globally Lipschitz in $(x,t)$ for $t\geq0$,
when defined by $u(x,0)=g(x)$.
\end{lemma}
\begin{expandedproof}
If $y$ minimizes at $x$, use $y+h$ as competitor at $x+h$.
The velocity is unchanged, so $u(x+h,t)-u(x,t)\leq K|h|$.
Exchange $x,x+h$ for the reverse inequality.

For time zero, the competitor $y=x$ gives
$u(x,t)\leq g(x)+tL(0)$, while duality gives
\[
 u(x,t)\geq g(x)+t\inf_v[L(v)-K|v|]
 =g(x)-t\max_{|p|\leq K}H(p).
\]
Thus $|u(x,t)-g(x)|\leq Ct$, uniformly in $x$, with
\[
 C=\max\{|L(0)|,\max_{|p|\leq K}|H(p)|\}.
\]
Apply the same argument to (DP), using the spatial Lipschitz bound
for $u(\cdot,s)$, to get $|u(x,t)-u(x,s)|\leq C|t-s|$.
\end{expandedproof}


\PDEBodyDiagram{3.3}{Equal minimizing costs can produce a corner in the value.}
{With $H(p)=p^2/2$, $t=1$ and $g(y)=\max(-|y|,-2)$, the Hopf--Lax cost at $x=0$ has two minimizers $y=\pm1$. Their candidate endpoint slopes $(x-y)/t$ differ. Indeed near $x=0$ the value is $u(x,1)=-1/2-|x|$, so it is not differentiable there.}
{ch03-hopflax-tie}

\begin{theorem}[PDE at differentiability points; source Theorems 5--6]
The Hopf--Lax function satisfies $u_t+H(Du)=0$ at every point where
it is differentiable, hence almost everywhere by Rademacher's theorem.
\end{theorem}
\begin{expandedproof}
At a differentiability point $(x,t)$, (DP) with the competitor $x$
over a short forward time gives
$u(x+hv,t+h)\leq hL(v)+u(x,t)$ for every fixed $v$.
Divide by $h>0$ and pass to zero:
$u_t+v\cdot Du\leq L(v)$. Supremizing in $v$ yields
$u_t+H(Du)\leq0$.

For the other sign choose a minimizing $y$ in (HL) at $(x,t)$ and
put $v_*=(x-y)/t$.
At time $t-h$, $z=x-hv_*$ lies on the straight minimizing path.
Using $y$ as competitor there yields
$u(x,t)-u(z,t-h)\geq hL(v_*)$.
Differentiability gives $u_t+v_*\cdot Du\geq L(v_*)$.
Since $H(Du)\geq v_*\cdot Du-L(v_*)$, the reverse sign follows.
Neither argument differentiates a potentially nonsmooth minimizing
point $y(x,t)$.
\end{expandedproof}
\Needspace{45mm}
\PDEAnnotation{最小值不一定可微}{多个最优起点可能同时出现，不能直接对 $y(x,t)$ 求导。动态规划用固定竞争者给一侧不等式，再沿一条最优路径给另一侧，避开了这个问题。}

For $H(p)=|p|^2/2$, $L(v)=|v|^2/2$ and
$u(x,t)=\min_y\{|x-y|^2/(2t)+g(y)\}$.
At a smooth minimizing branch,
$(x-y)/t=Dg(y)$, so $x=y+tDg(y)$ and $Du=Dg(y)$,
recovering the local characteristic formulas.
When branches compete, the minimum remains meaningful even though
their smooth projection has broken down.

\subsection{Semiconcavity and uniqueness (Evans \S3.3.3)}
\label{sec:3-semiconcavity}
An a.e. PDE alone is insufficient for uniqueness.
For $H(p)=p^2$ and $g=0$ in one dimension, both $u=0$ and
\[
 v(x,t)=\min\{0,|x|-t\}
\]
are Lipschitz and satisfy $v_t+v_x^2=0$ off the three lines
$x=0,\pm t$. At the central upward corner,
$v(h,t)+v(-h,t)-2v(0,t)=2|h|$ for $|h|<t$.
This cannot be bounded by a constant times $h^2$.

\begin{definition}
A function $f$ is semiconcave with constant $C$ if
\[
 f(x+z)-2f(x)+f(x-z)\leq C|z|^2 .
\]
For continuous $f$ this is equivalent to concavity of
$f(x)-C|x|^2/2$, by the midpoint criterion and continuity.
For smooth $f$ it is equivalent to $D^2f\leq CI$.
\end{definition}
\Needspace{45mm}
\PDEAnnotation{半凹只限制向上弯曲}{允许向下尖角，禁止二阶差分过大的向上尖角。它不是要求 Hessian 的绝对值有界，也不是要求函数本身凹；先减去一个二次函数后才凹。}

\begin{theorem}[Source Lemmas 3--4]
If $g$ is semiconcave with constant $C$, the Hopf--Lax function has
the same spatial constant for all $t$.
If instead $D^2H\geq\theta I$, $\theta>0$, it satisfies
\[
 u(x+z,t)-2u(x,t)+u(x-z,t)\leq\frac{|z|^2}{\theta t},
\]
even for initial $g$ that is merely Lipschitz.
\end{theorem}
\begin{expandedproof}[including the dual estimate of Problem 15]
For the first assertion choose a minimizer $y$ at $(x,t)$.
Use $y+z$ and $y-z$ as competitors at $x+z$ and $x-z$.
All three velocities are equal, so their $L$ terms cancel.
The remaining second difference is that of $g$ at $y$.

For the second assertion, uniform convexity makes the maximizer $p(v)$
in $L(v)=\max_p(v\cdot p-H(p))$ unique, with $v=DH(p(v))$.
Strong monotonicity gives
\[
 \theta|p(v)-p(w)|^2
 \leq (v-w)\cdot(p(v)-p(w)),
\]
so $|p(v)-p(w)|\leq|v-w|/\theta$.
Comparing the optimizing expressions at $v$ and $v+h$ shows
$DL(v)=p(v)$; equivalently, apply the implicit function theorem
to $DH(p)=v$. Hence $DL$ is $1/\theta$-Lipschitz.
Integrating its increment along a segment gives
\[
 L(v+h)\leq L(v)+DL(v)\cdot h+\frac{|h|^2}{2\theta}.
\]
Add the estimates for $h$ and $-h$:
$L(v+h)+L(v-h)-2L(v)\leq|h|^2/\theta$.
Now use the same minimizing $y$ at all three target points.
The initial terms cancel and $h=z/t$ gives the claimed constant.

In particular, taking $v=(v_1+v_2)/2$, $h=(v_1-v_2)/2$ proves
\[
 \tfrac12L(v_1)+\tfrac12L(v_2)
 \leq L((v_1+v_2)/2)+\frac{|v_1-v_2|^2}{8\theta}.
\]
This is the dual inequality omitted from the source lemma's proof.
\end{expandedproof}

\begin{theorem}[Uniqueness in the source weak class; source Theorems 7--8]
Let $H$ be smooth, convex and superlinear. Among globally Lipschitz
functions with the same initial datum, satisfying the HJ equation a.e.
and the spatial semiconcavity bound
\[
 u(x+z,t)-2u(x,t)+u(x-z,t)
 \leq C(1+t^{-1})|z|^2,
\]
there is at most one solution.
The Hopf--Lax function belongs to this class if $g$ is semiconcave
or $H$ is uniformly convex.
\end{theorem}
\begin{expandedproof}[the source's regularized cone argument]
Let $w=u-\widetilde u$. Subtract the a.e. PDE and integrate $DH$
on the gradient segment:
\[
 w_t+b\cdot Dw=0,\qquad
 b=\int_0^1 DH(\lambda Du+(1-\lambda)D\widetilde u)\,d\lambda .
\]
Both gradients lie in a fixed ball, so $|b|\leq R$ for a finite $R$.

\paragraph{1. Smooth the coefficients, preserving a one-sided bound.}
Locally convolve $u,\widetilde u$ in space-time and define $b_\varepsilon$
by the same expression with their smoothed spatial gradients.
For $t>2\varepsilon$, semiconcavity implies
\[
 D^2u_\varepsilon,\ D^2\widetilde u_\varepsilon
      \leq C'(1+t^{-1})I .
\]
To check this, integrate their second-difference inequalities against
the nonnegative smoothing kernel, bound $1/s$ by $2/t$ on its time
support, then divide by the squared increment and let it tend to zero.
The gradients remain bounded and converge a.e. to the original ones.

Writing $A=D^2H\geq0$ and $B$ for the convex combination of smoothed
Hessians gives $B\leq C'(1+t^{-1})I$.
Thus $\operatorname{tr}(AB)\leq C'(1+t^{-1})\operatorname{tr}A$,
by conjugating with $A^{1/2}$ and taking traces.
Since $A$ is bounded on the gradient ball,
$\operatorname{div}b_\varepsilon\leq C''(1+t^{-1})$.
Also $|b_\varepsilon|\leq R$ and $b_\varepsilon\to b$ a.e.

\Needspace{45mm}
\PDEAnnotation{凸性把单边 Hessian 界传给速度}{两个半凹函数的 Hessian 只需有上界。因 $D^2H$ 半正定，迹运算仍给出 $\operatorname{div}b$ 的上界；若失去凸性，这一步的符号便没有保证。}
\paragraph{2. Integrate a nonnegative transform in a shrinking cone.}
Take a smooth nonnegative function $\phi$, and put $v=\phi(w)$.
Weak Lipschitz calculus gives $v_t+b\cdot Dv=0$ a.e.
On $B_{\rho(t)}(x_0)$, $\rho(t)=r+R(T-t)$ with $r>0$, set
$e(t)=\int v$. The moving-boundary term and $|b_\varepsilon|\leq R$
make its boundary flux nonpositive. For a.e. $t$ this yields
\[
 e'(t)\leq C''(1+t^{-1})e(t)
       +\int_{B_{\rho(t)}}(b_\varepsilon-b)\cdot Dv.
\]
These integrations hold for Lipschitz $v$ by approximation, or its
Sobolev integration formula on the smooth ball.
For each fixed $\delta>0$, the error tends to zero in $L^1(\delta,T)$
by dominated convergence on the compact cone.
Integrate against a nonnegative time test function and let
$\varepsilon\to0$, obtaining the same differential inequality without
the error on $[\delta,T]$.

\paragraph{3. Start Gronwall away from the singular time zero.}
Let $M$ be the sum of the time Lipschitz constants.
Since $w(\cdot,0)=0$, $|w(x,\delta)|\leq M\delta$.
Choose $\phi=\phi_\delta$ zero on $[-M\delta,M\delta]$ and strictly
positive outside. Then $e(\delta)=0$.
Gronwall on $[\delta,T]$ has a finite coefficient integral, and gives
$e(T)=0$. Hence $|w|\leq M\delta$ on $B_r(x_0)$ at time $T$,
using continuity to pass from a.e. points.
Let $\delta\downarrow0$. The ball and $T$ were arbitrary, so $w=0$.
Existence in the stated class follows from the two semiconcavity
estimates and the preceding Hopf--Lax theorem.
\end{expandedproof}
\Needspace{45mm}
\PDEAnnotation{不能从零时刻硬套 Gronwall}{$1/t$ 在零附近不可积。证明先用时间 Lipschitz 界把差值截掉，在正时刻 $\delta$ 得到零初始能量，再应用 Gronwall，最后令 $\delta\to0$。}

This is the semiconcavity-based weak formulation used in this chapter.
Viscosity solutions provide a broader comparison framework in Chapter 10;
the a.e. equation and the present selection condition must not be
silently interchanged.

